The visualization and quantitative characterization of acoustic fields have
long been fundamental pursuits in experimental physics, engineering acoustics,
and fluid dynamics.
Because acoustic waves propagate as longitudinal pressure variations within a
compressible fluid medium, they are inherently invisible to direct optical
observation.
The endeavour to map these fields mathematically and visually traces its
origins to 1787, when the physicist Ernst Chladni introduced a technique
utilising finely dispersed sand sprinkled on vibrating metal plates.
As the plates were excited by the bow of a violin, the sand particles migrated
away from regions of maximum vibration (antinodes) and collected along nodal
lines---regions of zero displacement---providing the scientific community with
the first tangible, visual representation of acoustic resonance, modal
patterns, and structural vibrations~\cite{FernandezComesana2014}.
However, while Chladni's method was revolutionary for its time, it was limited
strictly to standing waves on solid boundaries and offered no mechanism for
mapping propagating acoustic fields in free space or fluid media.

